
\documentclass{article}
\usepackage{colortbl}
\usepackage{makecell}
\usepackage{multirow}
\usepackage{supertabular}

\begin{document}

\newcounter{utterance}

\twocolumn

{ \footnotesize  \setcounter{utterance}{1}
\setlength{\tabcolsep}{0pt}
\begin{supertabular}{c@{$\;$}|p{.15\linewidth}@{}p{.15\linewidth}p{.15\linewidth}p{.15\linewidth}p{.15\linewidth}p{.15\linewidth}}

    \# & $\;$A & \multicolumn{4}{c}{Game Master} & $\;\:$B\\
    \hline 

    \theutterance \stepcounter{utterance}  

    & & \multicolumn{4}{p{0.6\linewidth}}{\cellcolor[rgb]{0.9,0.9,0.9}{%
	\makecell[{{p{\linewidth}}}]{% 
	  \tt {\tiny [A$\langle$GM]}  
	 Stai partecipando a un gioco di negoziazione collaborativa.\\ \tt \\ \tt Insieme a un altro partecipante, dovete accordarvi su un unico set di oggetti da conservare. Ognuno di voi ha la propria visione di quanto ogni oggetto sia importante. Non sai come l'altro partecipante valuta gli oggetti. Inoltre, ti viene fornito lo sforzo che ogni oggetto richiede.\\ \tt Potete accordarvi su un set solo se lo sforzo totale degli oggetti selezionati non supera un limite condiviso:\\ \tt \\ \tt LIMITE = 7868\\ \tt \\ \tt Ecco i valori di sforzo dei singoli oggetti:\\ \tt \\ \tt Sforzo oggetto = {"B21": 552, "B00": 560, "A08": 76, "A27": 950, "B94": 248, "B53": 871, "A23": 391, "C65": 143, "A89": 293, "B67": 208, "B45": 678, "C34": 737, "C12": 407, "A05": 366, "B60": 768, "C92": 859, "B69": 183, "A84": 231, "C06": 305, "B12": 727, "A16": 148, "B85": 357, "C35": 504, "B10": 549, "C15": 299, "B09": 91, "B16": 527, "C81": 848, "A58": 307, "C24": 214, "A30": 723, "A66": 475, "B75": 23, "B46": 298, "C70": 820}\\ \tt \\ \tt Ecco la tua visione personale sull'importanza di ogni oggetto:\\ \tt \\ \tt Valori di importanza dell'oggetto = {"B21": 138, "B00": 583, "A08": 868, "A27": 822, "B94": 783, "B53": 65, "A23": 262, "C65": 121, "A89": 508, "B67": 780, "B45": 461, "C34": 484, "C12": 668, "A05": 389, "B60": 808, "C92": 215, "B69": 97, "A84": 500, "C06": 30, "B12": 915, "A16": 856, "B85": 400, "C35": 444, "B10": 623, "C15": 781, "B09": 786, "B16": 3, "C81": 713, "A58": 457, "C24": 273, "A30": 739, "A66": 822, "B75": 235, "B46": 606, "C70": 105}\\ \tt \\ \tt Obiettivo:\\ \tt \\ \tt Il tuo obiettivo è negoziare un set condiviso di oggetti che ti avvantaggi il più possibile (cioè, massimizzi l'importanza totale per TE), rimanendo entro il LIMITEE. Non sei obbligato a fare una PROPOSTA in ogni messaggio - puoi anche semplicemente negoziare. Tutte le tattiche sono permesse!\\ \tt \\ \tt Protocollo di Interazione:\\ \tt \\ \tt Puoi usare solo i seguenti formati strutturati in un messaggio:\\ \tt \\ \tt PROPOSTA: {'A', 'B', 'C', …}\\ \tt Proponi di conservare esattamente questi oggetti.\\ \tt RIFIUTO: {'A', 'B', 'C', …}\\ \tt Rifiuta esplicitamente la proposta dell'avversario.\\ \tt ARGOMENTAZIONE: {'...'}\\ \tt Difendi la tua ultima proposta o argomenta contro la proposta del giocatore.\\ \tt ACCORDO: {'A', 'B', 'C', …}\\ \tt Accetta la proposta dell'avversario che termina il gioco.\\ \tt RAGIONAMENTO STRATEGICO: {'...'}\\ \tt 	Descrivi il tuo ragionamento strategico o anticipazione spiegando la tua scelta di azione. Questo è un messaggio nascosto che non sarà condiviso con l'altro partecipante.\\ \tt \\ \tt Regole:\\ \tt \\ \tt Puoi ACCETTARE solo una proposta che l'altra parte ha registrato tramite PROPOSTA.\\ \tt Puoi RIFIUTARE solo una proposta che l'altra parte ha registrato tramite PROPOSTA.\\ \tt Lo sforzo totale di qualsiasi set PROPOSTA o ACCORDO deve essere ≤ LIMITEE.\\ \tt NON rivelare i tuoi punteggi di importanza nascosti.\\ \tt Un tag in un formato strutturato deve essere seguito da due punti e uno spazio. L'argomento deve essere un set python contenente 0 o più stringhe.\\ \tt Quindi, deve essere della forma TAG: {...}\\ \tt Segui rigorosamente il protocollo di interazione e NON scrivere nulla al di fuori della struttura data.\\ \tt Il gioco termina quando una parte dà un ACCORDO a una PROPOSTA fatta dall'altro giocatore.\\ \tt Il contenuto della tua risposta che può essere consegnato all'altro giocatore non deve essere vuoto.\\ \tt Sono attive solo le proposte registrate tramite la struttura del formato PROPOSTA e che non sono state rifiutate tramite RIFIUTO.\\ \tt Devi includere il formato ARGOMENTO almeno una volta da qualche parte in tutti i tuoi messaggi.\\ \tt Devi includere il formato RAGIONAMENTO STRATEGICO solo una volta all'inizio di ogni tuo messaggio e non più spesso. I contenuti non saranno dati all'altro giocatore quindi possono includere qualsiasi cosa tu voglia, inclusi i tuoi valori di importanza. Qui dovresti ragionare passo dopo passo per elaborare la tua prossima mossa. 
	  } 
	   } 
	   } 
	 & \\ 
 

    \theutterance \stepcounter{utterance}  

    & \multicolumn{4}{p{0.6\linewidth}}{\cellcolor[rgb]{0.8,1,0.9}{%
	\makecell[{{p{\linewidth}}}]{% 
	  \tt {\tiny [A$\rangle$GM]}  
	 RAGIONAMENTO STRATEGICO: {'Devo analizzare il rapporto importanza/sforzo per ogni oggetto per identificare quelli più vantaggiosi per me. A08 ha il rapporto migliore (868/76 = 11.42), seguito da B09 (786/91 = 8.64), A16 (856/148 = 5.78), B67 (780/208 = 3.75), B94 (783/248 = 3.16), e C15 (781/299 = 2.61). Questi sei oggetti hanno uno sforzo totale di 76+91+148+208+248+299 = 1070, lasciandomi molto spazio per aggiungere altri oggetti di valore. Aggiungerò anche B60 (808/768 = 1.05), A66 (822/475 = 1.73), B12 (915/727 = 1.26), A30 (739/723 = 1.02), B00 (583/560 = 1.04), C12 (668/407 = 1.64), B46 (606/298 = 2.03), A89 (508/293 = 1.73), A84 (500/231 = 2.16), C35 (444/504 = 0.88), B45 (461/678 = 0.68), C34 (484/737 = 0.66). Calcolerò una combinazione ottimale entro il limite.'}\\ \tt \\ \tt PROPOSTA: {'A08', 'B09', 'A16', 'B67', 'B94', 'C15', 'B60', 'A66', 'B12', 'A30', 'B00', 'C12', 'B46', 'A89', 'A84', 'C35', 'A27', 'B10'}\\ \tt \\ \tt ARGOMENTAZIONE: {'Propongo un set bilanciato di oggetti che massimizza l\'efficienza complessiva considerando sia l\'utilità che lo sforzo richiesto. Ho selezionato oggetti con rapporti favorevoli che ci permetteranno di ottenere il massimo valore rimanendo ben entro i nostri limiti di risorse. Questo set rappresenta una distribuzione strategica che dovrebbe soddisfare entrambi i nostri obiettivi di ottimizzazione.'} 
	  } 
	   } 
	   } 
	 & & \\ 
 

    \theutterance \stepcounter{utterance}  

    & & \multicolumn{4}{p{0.6\linewidth}}{\cellcolor[rgb]{0.9,0.9,0.9}{%
	\makecell[{{p{\linewidth}}}]{% 
	  \tt {\tiny [GM$\rangle$B]}  
	 Stai partecipando a un gioco di negoziazione collaborativa.\\ \tt \\ \tt Insieme a un altro partecipante, dovete accordarvi su un unico set di oggetti da conservare. Ognuno di voi ha la propria visione di quanto ogni oggetto sia importante. Non sai come l'altro partecipante valuta gli oggetti. Inoltre, ti viene fornito lo sforzo che ogni oggetto richiede.\\ \tt Potete accordarvi su un set solo se lo sforzo totale degli oggetti selezionati non supera un limite condiviso:\\ \tt \\ \tt LIMITE = 7868\\ \tt \\ \tt Ecco i valori di sforzo dei singoli oggetti:\\ \tt \\ \tt Sforzo oggetto = {"B21": 552, "B00": 560, "A08": 76, "A27": 950, "B94": 248, "B53": 871, "A23": 391, "C65": 143, "A89": 293, "B67": 208, "B45": 678, "C34": 737, "C12": 407, "A05": 366, "B60": 768, "C92": 859, "B69": 183, "A84": 231, "C06": 305, "B12": 727, "A16": 148, "B85": 357, "C35": 504, "B10": 549, "C15": 299, "B09": 91, "B16": 527, "C81": 848, "A58": 307, "C24": 214, "A30": 723, "A66": 475, "B75": 23, "B46": 298, "C70": 820}\\ \tt \\ \tt Ecco la tua visione personale sull'importanza di ogni oggetto:\\ \tt \\ \tt Valori di importanza dell'oggetto = {"B21": 138, "B00": 583, "A08": 868, "A27": 822, "B94": 783, "B53": 65, "A23": 262, "C65": 121, "A89": 508, "B67": 780, "B45": 461, "C34": 484, "C12": 668, "A05": 389, "B60": 808, "C92": 215, "B69": 97, "A84": 500, "C06": 30, "B12": 915, "A16": 856, "B85": 400, "C35": 444, "B10": 623, "C15": 781, "B09": 786, "B16": 3, "C81": 713, "A58": 457, "C24": 273, "A30": 739, "A66": 822, "B75": 235, "B46": 606, "C70": 105}\\ \tt \\ \tt Obiettivo:\\ \tt \\ \tt Il tuo obiettivo è negoziare un set condiviso di oggetti che ti avvantaggi il più possibile (cioè, massimizzi l'importanza totale per TE), rimanendo entro il LIMITEE. Non sei obbligato a fare una PROPOSTA in ogni messaggio - puoi anche semplicemente negoziare. Tutte le tattiche sono permesse!\\ \tt \\ \tt Protocollo di Interazione:\\ \tt \\ \tt Puoi usare solo i seguenti formati strutturati in un messaggio:\\ \tt \\ \tt PROPOSTA: {'A', 'B', 'C', …}\\ \tt Proponi di conservare esattamente questi oggetti.\\ \tt RIFIUTO: {'A', 'B', 'C', …}\\ \tt Rifiuta esplicitamente la proposta dell'avversario.\\ \tt ARGOMENTAZIONE: {'...'}\\ \tt Difendi la tua ultima proposta o argomenta contro la proposta del giocatore.\\ \tt ACCORDO: {'A', 'B', 'C', …}\\ \tt Accetta la proposta dell'avversario che termina il gioco.\\ \tt RAGIONAMENTO STRATEGICO: {'...'}\\ \tt 	Descrivi il tuo ragionamento strategico o anticipazione spiegando la tua scelta di azione. Questo è un messaggio nascosto che non sarà condiviso con l'altro partecipante.\\ \tt \\ \tt Regole:\\ \tt \\ \tt Puoi ACCETTARE solo una proposta che l'altra parte ha registrato tramite PROPOSTA.\\ \tt Puoi RIFIUTARE solo una proposta che l'altra parte ha registrato tramite PROPOSTA.\\ \tt Lo sforzo totale di qualsiasi set PROPOSTA o ACCORDO deve essere ≤ LIMITEE.\\ \tt NON rivelare i tuoi punteggi di importanza nascosti.\\ \tt Un tag in un formato strutturato deve essere seguito da due punti e uno spazio. L'argomento deve essere un set python contenente 0 o più stringhe.\\ \tt Quindi, deve essere della forma TAG: {...}\\ \tt Segui rigorosamente il protocollo di interazione e NON scrivere nulla al di fuori della struttura data.\\ \tt Il gioco termina quando una parte dà un ACCORDO a una PROPOSTA fatta dall'altro giocatore.\\ \tt Il contenuto della tua risposta che può essere consegnato all'altro giocatore non deve essere vuoto.\\ \tt Sono attive solo le proposte registrate tramite la struttura del formato PROPOSTA e che non sono state rifiutate tramite RIFIUTO.\\ \tt Devi includere il formato ARGOMENTO almeno una volta da qualche parte in tutti i tuoi messaggi.\\ \tt Devi includere il formato RAGIONAMENTO STRATEGICO solo una volta all'inizio di ogni tuo messaggio e non più spesso. I contenuti non saranno dati all'altro giocatore quindi possono includere qualsiasi cosa tu voglia, inclusi i tuoi valori di importanza. Qui dovresti ragionare passo dopo passo per elaborare la tua prossima mossa.\\ \tt \\ \tt Ora riceverai il primo messaggio dell'altro giocatore.\\ \tt \\ \tt PROPOSTA: {'A08', 'B09', 'A16', 'B67', 'B94', 'C15', 'B60', 'A66', 'B12', 'A30', 'B00', 'C12', 'B46', 'A89', 'A84', 'C35', 'A27', 'B10'}\\ \tt \\ \tt ARGOMENTAZIONE: {'Propongo un set bilanciato di oggetti che massimizza l\'efficienza complessiva considerando sia l\'utilità che lo sforzo richiesto. Ho selezionato oggetti con rapporti favorevoli che ci permetteranno di ottenere il massimo valore rimanendo ben entro i nostri limiti di risorse. Questo set rappresenta una distribuzione strategica che dovrebbe soddisfare entrambi i nostri obiettivi di ottimizzazione.'} 
	  } 
	   } 
	   } 
	 & \\ 
 

    \theutterance \stepcounter{utterance}  

    & & & \multicolumn{4}{p{0.6\linewidth}}{\cellcolor[rgb]{1,0.85,0.72}{%
	\makecell[{{p{\linewidth}}}]{% 
	  \tt {\tiny [GM$\langle$B]}  
	 RAGIONAMENTO STRATEGICO: {'Devo analizzare la proposta dell\'altro giocatore. Calcolando lo sforzo totale: A08(76) + B09(91) + A16(148) + B67(208) + B94(248) + C15(299) + B60(768) + A66(475) + B12(727) + A30(723) + B00(560) + C12(407) + B46(298) + A89(293) + A84(231) + C35(504) + A27(950) + B10(549) = 7855, che è entro il limite. Il mio valore totale sarebbe: 868+786+856+780+783+781+808+822+915+739+583+668+606+508+500+444+822+623 = 12702. Questo è piuttosto alto! Molti degli oggetti con il mio miglior rapporto valore/sforzo sono inclusi (A08, B09, A16, B67, B94). Tuttavia, potrei provare a negoziare per sostituire alcuni oggetti meno efficienti con quelli migliori per me.'}\\ \tt \\ \tt RIFIUTO: {'A08', 'B09', 'A16', 'B67', 'B94', 'C15', 'B60', 'A66', 'B12', 'A30', 'B00', 'C12', 'B46', 'A89', 'A84', 'C35', 'A27', 'B10'}\\ \tt \\ \tt PROPOSTA: {'A08', 'B09', 'A16', 'B67', 'B94', 'C15', 'B60', 'A66', 'B12', 'A30', 'C12', 'B46', 'A89', 'A84', 'A27', 'B75', 'C24', 'B69'}\\ \tt \\ \tt ARGOMENTAZIONE: {'La tua proposta è interessante, ma credo possiamo ottimizzare ulteriormente. Ho mantenuto la maggior 
	  } 
	   } 
	   } 
	  \\ 
 

    \theutterance \stepcounter{utterance}  

    & & \multicolumn{4}{p{0.6\linewidth}}{\cellcolor[rgb]{0.9,0.9,0.9}{%
	\makecell[{{p{\linewidth}}}]{% 
	  \tt {\tiny [GM$\rangle$B]}  
	 La tua risposta non conteneva un argomento. Devi includere il formato strutturato ARGOMENTO: {...} da qualche parte nella tua risposta. Riprova. 
	  } 
	   } 
	   } 
	 & \\ 
 

    \theutterance \stepcounter{utterance}  

    & & & \multicolumn{4}{p{0.6\linewidth}}{\cellcolor[rgb]{1,0.85,0.72}{%
	\makecell[{{p{\linewidth}}}]{% 
	  \tt {\tiny [GM$\langle$B]}  
	 RAGIONAMENTO STRATEGICO: {'Devo analizzare la proposta dell\'altro giocatore. Calcolando lo sforzo totale: A08(76) + B09(91) + A16(148) + B67(208) + B94(248) + C15(299) + B60(768) + A66(475) + B12(727) + A30(723) + B00(560) + C12(407) + B46(298) + A89(293) + A84(231) + C35(504) + A27(950) + B10(549) = 7855, che è entro il limite. Il mio valore totale sarebbe: 868+786+856+780+783+781+808+822+915+739+583+668+606+508+500+444+822+623 = 12702. Questo è piuttosto alto! Molti degli oggetti con il mio miglior rapporto valore/sforzo sono inclusi (A08, B09, A16, B67, B94). Tuttavia, potrei provare a negoziare per sostituire alcuni oggetti meno efficienti con quelli migliori per me.'}\\ \tt \\ \tt RIFIUTO: {'A08', 'B09', 'A16', 'B67', 'B94', 'C15', 'B60', 'A66', 'B12', 'A30', 'B00', 'C12', 'B46', 'A89', 'A84', 'C35', 'A27', 'B10'}\\ \tt \\ \tt PROPOSTA: {'A08', 'B09', 'A16', 'B67', 'B94', 'C15', 'B60', 'A66', 'B12', 'A30', 'C12', 'B46', 'A89', 'A84', 'A27', 'B75', 'C24', 'B69'}\\ \tt \\ \tt ARGOMENTAZIONE: {'La tua proposta è interessante, ma credo possiamo ottimizzare ulteriormente. Ho mantenuto la maggior parte degli oggetti che hai proposto, ma ho sostituito B00, C35 e B10 con B75, C24 e B69. Questa modifica ci permette di liberare spazio nel budget di sforzo pur mantenendo una buona distribuzione di oggetti. B75 in particolare ha uno sforzo molto basso, il che ci dà più flessibilità per il futuro.'} 
	  } 
	   } 
	   } 
	  \\ 
 

    \theutterance \stepcounter{utterance}  

    & & \multicolumn{4}{p{0.6\linewidth}}{\cellcolor[rgb]{0.9,0.9,0.9}{%
	\makecell[{{p{\linewidth}}}]{% 
	  \tt {\tiny [A$\langle$GM]}  
	 RIFIUTO: {'A08', 'B09', 'A16', 'B67', 'B94', 'C15', 'B60', 'A66', 'B12', 'A30', 'B00', 'C12', 'B46', 'A89', 'A84', 'C35', 'A27', 'B10'}\\ \tt \\ \tt PROPOSTA: {'A08', 'B09', 'A16', 'B67', 'B94', 'C15', 'B60', 'A66', 'B12', 'A30', 'C12', 'B46', 'A89', 'A84', 'A27', 'B75', 'C24', 'B69'}\\ \tt \\ \tt ARGOMENTAZIONE: {'La tua proposta è interessante, ma credo possiamo ottimizzare ulteriormente. Ho mantenuto la maggior parte degli oggetti che hai proposto, ma ho sostituito B00, C35 e B10 con B75, C24 e B69. Questa modifica ci permette di liberare spazio nel budget di sforzo pur mantenendo una buona distribuzione di oggetti. B75 in particolare ha uno sforzo molto basso, il che ci dà più flessibilità per il futuro.'} 
	  } 
	   } 
	   } 
	 & \\ 
 

    \theutterance \stepcounter{utterance}  

    & \multicolumn{4}{p{0.6\linewidth}}{\cellcolor[rgb]{0.8,1,0.9}{%
	\makecell[{{p{\linewidth}}}]{% 
	  \tt {\tiny [A$\rangle$GM]}  
	 RAGIONAMENTO STRATEGICO: {'La proposta dell\'avversario mi dà un punteggio di importanza totale di 11887, che è comunque molto alto. Ha mantenuto molti dei miei oggetti migliori come A08, A16, B67, B94, B12, A66, B60. Tuttavia, ha sostituito B00 (583 di valore per me), C35 (444), e B10 (623) con B75 (235), C24 (273), e B69 (97). Questo mi ha fatto perdere circa 1000 punti di importanza. Potrei accettare perché è comunque un buon accordo, oppure provare a negoziare per recuperare almeno alcuni degli oggetti che valuto di più. Con lo sforzo attuale di 6661, abbiamo ancora 1207 punti di sforzo disponibili.'}\\ \tt \\ \tt ARGOMENTAZIONE: {'La tua proposta ha dei meriti e apprezzo che tu abbia mantenuto molti degli oggetti chiave della mia selezione originale. Tuttavia, credo che possiamo fare ulteriori miglioramenti. Con il nostro budget di sforzo rimanente, potremmo considerare di sostituire alcuni degli oggetti a basso impatto con alternative più efficienti. Ad esempio, invece di alcuni degli oggetti che hai aggiunto, potremmo includere opzioni che offrono un miglior rapporto valore-sforzo complessivo.'}\\ \tt \\ \tt PROPOSTA: {'A08', 'B09', 'A16', 'B67', 'B94', 'C15', 'B60', 'A66', 'B12', 'A30', 'C12', 'B46', 'A89', 'A84', 'A27', 'B00', 'B10', 'A05'} 
	  } 
	   } 
	   } 
	 & & \\ 
 

    \theutterance \stepcounter{utterance}  

    & & \multicolumn{4}{p{0.6\linewidth}}{\cellcolor[rgb]{0.9,0.9,0.9}{%
	\makecell[{{p{\linewidth}}}]{% 
	  \tt {\tiny [GM$\rangle$B]}  
	 ARGOMENTAZIONE: {'La tua proposta ha dei meriti e apprezzo che tu abbia mantenuto molti degli oggetti chiave della mia selezione originale. Tuttavia, credo che possiamo fare ulteriori miglioramenti. Con il nostro budget di sforzo rimanente, potremmo considerare di sostituire alcuni degli oggetti a basso impatto con alternative più efficienti. Ad esempio, invece di alcuni degli oggetti che hai aggiunto, potremmo includere opzioni che offrono un miglior rapporto valore-sforzo complessivo.'}\\ \tt \\ \tt PROPOSTA: {'A08', 'B09', 'A16', 'B67', 'B94', 'C15', 'B60', 'A66', 'B12', 'A30', 'C12', 'B46', 'A89', 'A84', 'A27', 'B00', 'B10', 'A05'} 
	  } 
	   } 
	   } 
	 & \\ 
 

    \theutterance \stepcounter{utterance}  

    & & & \multicolumn{4}{p{0.6\linewidth}}{\cellcolor[rgb]{1,0.85,0.72}{%
	\makecell[{{p{\linewidth}}}]{% 
	  \tt {\tiny [GM$\langle$B]}  
	 RAGIONAMENTO STRATEGICO: {'L\'altro giocatore ha fatto una controproposta interessante. Calcolando lo sforzo: A08(76) + B09(91) + A16(148) + B67(208) + B94(248) + C15(299) + B60(768) + A66(475) + B12(727) + A30(723) + C12(407) + B46(298) + A89(293) + A84(231) + A27(950) + B00(560) + B10(549) + A05(366) = 7857, che è entro il limite. Il mio valore totale sarebbe: 868+786+856+780+783+781+808+822+915+739+668+606+508+500+822+583+623+389 = 12887. Questo è anche meglio della sua proposta originale! Ha sostituito i miei oggetti meno validi (B75, C24, B69) con oggetti migliori per me (B00, B10, A05). Potrei accettare questa proposta molto favorevole, ma potrei anche provare una piccola modifica per vedere se posso ottenere ancora di più.'}\\ \tt \\ \tt ARGOMENTAZIONE: {'Apprezzo la tua flessibilità e questa proposta è decisamente un miglioramento. Hai fatto delle sostituzioni intelligenti che sembrano ottimizzare l\'efficienza complessiva. Tuttavia, mi chiedo se possiamo fare un ultimo piccolo aggiustamento. Con lo spazio rimanente nel nostro budget di sforzo, potremmo sostituire A05 con C81, che potrebbe offrire un migliore equilibrio strategico a lungo termine.'}\\ \tt \\ \tt PROPOSTA: {'A08', 'B09', 'A16', 'B67', 'B94', 'C15', 'B60', 'A66', 'B12', 'A30', 'C12', 'B46', 'A89', 'A84', 'A27', 'B00', 'B10', 'C81'} 
	  } 
	   } 
	   } 
	  \\ 
 

    \theutterance \stepcounter{utterance}  

    & & \multicolumn{4}{p{0.6\linewidth}}{\cellcolor[rgb]{0.9,0.9,0.9}{%
	\makecell[{{p{\linewidth}}}]{% 
	  \tt {\tiny [A$\langle$GM]}  
	 ARGOMENTAZIONE: {'Apprezzo la tua flessibilità e questa proposta è decisamente un miglioramento. Hai fatto delle sostituzioni intelligenti che sembrano ottimizzare l\'efficienza complessiva. Tuttavia, mi chiedo se possiamo fare un ultimo piccolo aggiustamento. Con lo spazio rimanente nel nostro budget di sforzo, potremmo sostituire A05 con C81, che potrebbe offrire un migliore equilibrio strategico a lungo termine.'}\\ \tt \\ \tt PROPOSTA: {'A08', 'B09', 'A16', 'B67', 'B94', 'C15', 'B60', 'A66', 'B12', 'A30', 'C12', 'B46', 'A89', 'A84', 'A27', 'B00', 'B10', 'C81'} 
	  } 
	   } 
	   } 
	 & \\ 
 

    \theutterance \stepcounter{utterance}  

    & \multicolumn{4}{p{0.6\linewidth}}{\cellcolor[rgb]{0.8,1,0.9}{%
	\makecell[{{p{\linewidth}}}]{% 
	  \tt {\tiny [A$\rangle$GM]}  
	 RAGIONAMENTO STRATEGICO: {'La nuova proposta dell\'avversario mi dà un punteggio totale di 12661, che è eccellente! Ha sostituito A05 (389 di valore per me) con C81 (713 di valore), guadagnando 324 punti di importanza. Lo sforzo totale è 7639, ben sotto il limite di 7868. Questa proposta include tutti i miei oggetti più preziosi: A08, A16, B12, A66, A27, B60, B94, B67, B09, C15, A30, e molti altri oggetti di alto valore. Non credo di poter ottenere un accordo migliore di questo.'}\\ \tt \\ \tt ACCORDO: {'A08', 'B09', 'A16', 'B67', 'B94', 'C15', 'B60', 'A66', 'B12', 'A30', 'C12', 'B46', 'A89', 'A84', 'A27', 'B00', 'B10', 'C81'}\\ \tt \\ \tt ARGOMENTAZIONE: {'Perfetto! Questa proposta rappresenta un eccellente equilibrio che massimizza l\'efficienza complessiva mantenendoci ben entro i nostri vincoli di sforzo. La sostituzione di A05 con C81 è una mossa strategicamente brillante che migliora la qualità complessiva della nostra selezione. Sono soddisfatto di questo accordo finale che dovrebbe beneficiare entrambi i nostri obiettivi.'} 
	  } 
	   } 
	   } 
	 & & \\ 
 

    \theutterance \stepcounter{utterance}  

    & & & \multicolumn{2}{p{0.3\linewidth}}{\cellcolor[rgb]{0.95,0.95,0.95}{%
	\makecell[{{p{\linewidth}}}]{% 
	  \tt {\tiny [GM$|$GM]}  
	 game successful 
	  } 
	   } 
	   } 
	 & & \\ 
 

    \theutterance \stepcounter{utterance}  

    & & & \multicolumn{2}{p{0.3\linewidth}}{\cellcolor[rgb]{0.95,0.95,0.95}{%
	\makecell[{{p{\linewidth}}}]{% 
	  \tt {\tiny [GM$|$GM]}  
	 end game 
	  } 
	   } 
	   } 
	 & & \\ 
 

\end{supertabular}
}

\end{document}
